\documentclass[12pt, a4paper]{article}

\usepackage[utf8]{inputenc}
\usepackage[T1]{fontenc}
\usepackage[french]{babel}
\usepackage[margin=2.5cm]{geometry}
\usepackage{xcolor}
\usepackage{booktabs}
\usepackage{longtable}
\usepackage{array}
\usepackage{enumitem}
\usepackage{titlesec}
\usepackage{fancyhdr}
\usepackage{mdframed}
\usepackage{tikz}
\usetikzlibrary{shapes.geometric, arrows.meta, positioning, fit, backgrounds}

% ── Couleurs ──────────────────────────────────────────────────────────────────
\definecolor{dxcpurple}{RGB}{102, 0, 153}
\definecolor{blockbg}{RGB}{248, 242, 255}
\definecolor{greenbg}{RGB}{232, 245, 233}
\definecolor{orangebg}{RGB}{255, 243, 224}
\definecolor{redbg}{RGB}{255, 235, 235}
\definecolor{greytext}{RGB}{90, 90, 90}
\definecolor{greenborder}{RGB}{56, 142, 60}
\definecolor{orangeborder}{RGB}{230, 81, 0}
\definecolor{redborder}{RGB}{198, 40, 40}

% ── En-tête / pied de page ────────────────────────────────────────────────────
\pagestyle{fancy}
\fancyhf{}
\fancyhead[L]{\textcolor{dxcpurple}{\textbf{DXC Technology}}}
\fancyhead[R]{\textcolor{dxcpurple}{Architecture du Système de Reporting}}
\fancyfoot[C]{\thepage}
\renewcommand{\headrulewidth}{0.4pt}
\renewcommand{\headrule}{%
  \hbox to\headwidth{\color{dxcpurple}\leaders\hrule height \headrulewidth\hfill}}

% ── Titres ────────────────────────────────────────────────────────────────────
\titleformat{\section}{\large\bfseries\color{dxcpurple}}{}{0em}{}[\titlerule]
\titleformat{\subsection}{\normalsize\bfseries\color{dxcpurple}}{}{0em}{}

% ── Environnements encadrés ───────────────────────────────────────────────────
\newmdenv[
  backgroundcolor=blockbg,
  linecolor=dxcpurple,
  linewidth=1.2pt,
  innerleftmargin=12pt, innerrightmargin=12pt,
  innertopmargin=10pt, innerbottommargin=10pt,
  roundcorner=4pt
]{purplebox}

\newmdenv[
  backgroundcolor=greenbg,
  linecolor=greenborder,
  linewidth=1pt,
  innerleftmargin=12pt, innerrightmargin=12pt,
  innertopmargin=8pt, innerbottommargin=8pt,
  roundcorner=4pt
]{greenbox}

\newmdenv[
  backgroundcolor=orangebg,
  linecolor=orangeborder,
  linewidth=1pt,
  innerleftmargin=12pt, innerrightmargin=12pt,
  innertopmargin=8pt, innerbottommargin=8pt,
  roundcorner=4pt
]{orangebox}

\newmdenv[
  backgroundcolor=redbg,
  linecolor=redborder,
  linewidth=1pt,
  innerleftmargin=12pt, innerrightmargin=12pt,
  innertopmargin=8pt, innerbottommargin=8pt,
  roundcorner=4pt
]{redbox}

% ─────────────────────────────────────────────────────────────────────────────
\begin{document}

% ── Page de titre ─────────────────────────────────────────────────────────────
\begin{center}
  {\color{dxcpurple}\rule{\linewidth}{1.5pt}}\\[0.5cm]
  {\LARGE\bfseries\color{dxcpurple}
    Choix d'Architecture du Système de Reporting}\\[0.25cm]
  {\large IRPAUTO · DXC GraphTalk Service Desk}\\[0.15cm]
  {\normalsize\color{greytext}
    Analyse comparative des architectures envisagées et justification du choix final}\\[0.4cm]
  {\color{dxcpurple}\rule{\linewidth}{1.5pt}}
\end{center}

\vspace{0.6cm}

% ─────────────────────────────────────────────────────────────────────────────
\section{Introduction}

Dans le cadre du projet de tableau de bord de pilotage de sprint, plusieurs architectures
techniques ont été étudiées pour répondre aux besoins de collecte, de stockage et de
visualisation des données issues de Jira. Ce document présente les trois architectures
envisagées, leurs avantages et inconvénients respectifs, ainsi que la justification du
choix final retenu.

\vspace{0.3cm}

Les critères d'évaluation retenus sont les suivants~:

\begin{itemize}[leftmargin=2em, itemsep=3pt]
  \item \textbf{Sécurité et confidentialité des données}
  \item \textbf{Faisabilité technique} dans le contexte d'un stage
  \item \textbf{Coût} de mise en \oe uvre et d'exploitation
  \item \textbf{Performance} et capacité de montée en charge
  \item \textbf{Maintenabilité} et évolutivité de la solution
  \item \textbf{Disponibilité} du tableau de bord pour les parties prenantes
\end{itemize}

% ─────────────────────────────────────────────────────────────────────────────
\section{Architecture 1 — Base de données cloud externe (Aiven)}

\subsection{Description}

Dans cette architecture, les données extraites de Jira via l'API REST sont traitées
par un pipeline Python (ETL) et stockées dans une base de données MySQL hébergée sur
\textbf{Aiven}, une plateforme cloud managée. Microsoft Power BI Service se connecte
directement à cette base via Internet pour actualiser les rapports.

\vspace{0.2cm}

\begin{center}
\begin{tikzpicture}[
  node distance=1.2cm,
  box/.style={rectangle, rounded corners=4pt, draw=dxcpurple, fill=blockbg,
              minimum width=2.8cm, minimum height=0.9cm, align=center,
              font=\small},
  cloud/.style={rectangle, rounded corners=4pt, draw=orangeborder, fill=orangebg,
               minimum width=2.8cm, minimum height=0.9cm, align=center,
               font=\small},
  arr/.style={-{Stealth[length=5pt]}, thick, color=dxcpurple}
]
  \node[box] (jira)   {Jira REST API};
  \node[box, right=of jira] (python) {Python ETL};
  \node[cloud, right=of python] (aiven) {MySQL\\Aiven (cloud)};
  \node[box, right=of aiven] (pbi)   {Power BI\\Service};

  \draw[arr] (jira)   -- (python);
  \draw[arr] (python) -- (aiven);
  \draw[arr] (aiven)  -- (pbi);
\end{tikzpicture}
\end{center}

\vspace{0.2cm}

\subsection{Avantages et inconvénients}

\begin{orangebox}
\begin{itemize}[leftmargin=1.5em, itemsep=3pt]
  \item[\textcolor{greenborder}{+}] \textbf{Disponibilité permanente}~: la base est
        accessible 24h/24, Power BI Service peut se rafraîchir automatiquement sans
        intervention humaine.
  \item[\textcolor{greenborder}{+}] \textbf{Simplicité de connexion}~: Power BI
        Desktop et Service se connectent directement via l'IP publique, sans
        configuration réseau complexe.
  \item[\textcolor{greenborder}{+}] \textbf{Tier gratuit disponible}~: Aiven propose
        un tier gratuit suffisant pour le volume de données du projet.
  \item[\textcolor{redborder}{--}] \textbf{Risque de conformité}~: les données Jira
        contiennent des noms d'employés (données personnelles RGPD) et des références
        clients. Leur hébergement chez un prestataire cloud tiers nécessite une
        autorisation explicite de DXC Technology.
  \item[\textcolor{redborder}{--}] \textbf{Dépendance externe}~: en cas de panne
        d'Aiven ou de résiliation du service, l'ensemble du système est impacté.
  \item[\textcolor{redborder}{--}] \textbf{Token Jira et credentials exposés}~: les
        identifiants de connexion doivent être configurés dans un environnement externe.
\end{itemize}
\end{orangebox}

% ─────────────────────────────────────────────────────────────────────────────
\section{Architecture 2 — Power BI natif avec connecteurs directs}

\subsection{Description}

Dans cette architecture, Power BI se connecte directement à Jira via un connecteur
natif ou un connecteur tiers (Exalate, Alpha Serve). Toute la logique de transformation
et de calcul est réalisée dans Power Query (langage M) et DAX, sans pipeline Python
ni base de données intermédiaire.

\vspace{0.2cm}

\begin{center}
\begin{tikzpicture}[
  node distance=1.2cm,
  box/.style={rectangle, rounded corners=4pt, draw=dxcpurple, fill=blockbg,
              minimum width=2.8cm, minimum height=0.9cm, align=center,
              font=\small},
  arr/.style={-{Stealth[length=5pt]}, thick, color=dxcpurple}
]
  \node[box] (jira)  {Jira REST API};
  \node[box, right=2cm of jira] (pq) {Power Query\\(M + DAX)};
  \node[box, right=2cm of pq]   (pbi) {Power BI\\Service};

  \draw[arr] (jira) -- node[above, font=\tiny]{connecteur} (pq);
  \draw[arr] (pq)   -- (pbi);
\end{tikzpicture}
\end{center}

\vspace{0.2cm}

\subsection{Avantages et inconvénients}

\begin{orangebox}
\begin{itemize}[leftmargin=1.5em, itemsep=3pt]
  \item[\textcolor{greenborder}{+}] \textbf{Pas d'infrastructure}~: aucun serveur
        ni base de données à gérer.
  \item[\textcolor{greenborder}{+}] \textbf{Outil unique}~: toute la chaîne tient
        dans Power BI.
  \item[\textcolor{redborder}{--}] \textbf{Volume de données incompatible}~: le
        projet IRPAUTO contient plus de 55\,000 tickets. Power Query atteint ses
        limites au-delà de 10\,000--15\,000 lignes en mode Direct Query, entraînant
        des timeouts et des lenteurs.
  \item[\textcolor{redborder}{--}] \textbf{Logique métier non faisable}~: le matching
        flou des noms, le calcul du burndown avec exclusions, l'affectation des squads
        — toute cette logique est triviale en Python mais quasi impossible à maintenir
        en DAX/M.
  \item[\textcolor{redborder}{--}] \textbf{Coût}~: les connecteurs Jira tiers sont
        payants (\$10--50/mois). Power BI Premium est requis pour certaines
        fonctionnalités avancées.
  \item[\textcolor{redborder}{--}] \textbf{Non maintenable}~: le code M est difficile
        à lire, à déboguer et à versionner.
  \item[\textcolor{redborder}{--}] \textbf{Données exposées}~: les données transitent
        directement vers les serveurs Microsoft sans contrôle intermédiaire.
\end{itemize}
\end{orangebox}

% ─────────────────────────────────────────────────────────────────────────────
\section{Architecture 3 — SharePoint comme base de données}

\subsection{Description}

Microsoft SharePoint, déjà licencié par DXC, propose des \emph{Listes} qui peuvent
servir de stockage tabulaire. L'idée était d'y écrire les données transformées depuis
Python et de connecter Power BI aux listes SharePoint.

\subsection{Pourquoi cette architecture a été écartée}

\begin{redbox}
\begin{itemize}[leftmargin=1.5em, itemsep=3pt]
  \item \textbf{Limite de 5\,000 lignes}~: SharePoint Lists impose un seuil de vue
        à 5\,000 éléments. Au-delà, les requêtes échouent ou deviennent extrêmement
        lentes. Incompatible avec 55\,000 tickets.
  \item \textbf{Pas un vrai SGBD}~: SharePoint n'est pas une base de données
        relationnelle. Il ne supporte pas SQL, les types de données sont peu stricts
        et les écritures concurrentes peuvent corrompre les données.
  \item \textbf{Connexion Power BI lente}~: la connexion Power BI -- SharePoint
        passe par une API REST, bien plus lente qu'une connexion MySQL directe.
  \item \textbf{Écriture Python complexe}~: écrire dans SharePoint depuis Python
        nécessite des appels API Microsoft Graph avec une authentification OAuth
        complexe, contre une simple ligne \texttt{pandas.to\_sql()} pour MySQL.
\end{itemize}
\end{redbox}

% ─────────────────────────────────────────────────────────────────────────────
\section{Architecture Retenue — MySQL local via Docker}

\subsection{Description}

L'architecture finale retenue repose sur un conteneur \textbf{Docker} hébergeant
une instance \textbf{MySQL 8.0} sur un poste de travail dédié au sein du réseau
DXC. Power BI Desktop et Power BI Service se connectent à cette base via un
\textbf{On-premises Data Gateway} installé sur le même poste.

\vspace{0.3cm}

\begin{purplebox}
\textbf{Principe fondamental~:} aucune donnée ne quitte le réseau interne de
DXC Technology. Le pipeline complet — collecte, transformation, stockage,
visualisation — est maîtrisé de bout en bout par l'équipe.
\end{purplebox}

\vspace{0.4cm}

\subsection{Schéma d'architecture}

\begin{center}
\begin{tikzpicture}[
  node distance=1.0cm and 1.4cm,
  box/.style={rectangle, rounded corners=4pt, draw=dxcpurple, fill=blockbg,
              minimum width=3.0cm, minimum height=0.9cm, align=center, font=\small},
  intbox/.style={rectangle, rounded corners=4pt, draw=greenborder, fill=greenbg,
                 minimum width=3.0cm, minimum height=0.9cm, align=center, font=\small},
  cloudbox/.style={rectangle, rounded corners=4pt, draw=orangeborder, fill=orangebg,
                   minimum width=3.0cm, minimum height=0.9cm, align=center, font=\small},
  arr/.style={-{Stealth[length=5pt]}, thick, color=dxcpurple},
  garr/.style={-{Stealth[length=5pt]}, thick, color=greenborder, dashed}
]
  % Jira
  \node[box] (jira) {Jira REST API\\(Atlassian cloud)};

  % Python ETL
  \node[intbox, below=of jira] (etl) {Python ETL\\(collecte + calculs)};

  % Docker MySQL
  \node[intbox, right=2.2cm of etl] (mysql) {MySQL 8.0\\(Docker container)};

  % Streamlit
  \node[intbox, below=of etl] (streamlit) {Streamlit\\(tableau de bord)};

  % Gateway
  \node[intbox, below=of mysql] (gateway) {On-premises\\Data Gateway};

  % PowerBI Desktop
  \node[box, right=2.2cm of mysql] (pbd) {Power BI\\Desktop};

  % PowerBI Service
  \node[cloudbox, below=of pbd] (pbs) {Power BI\\Service (cloud)};

  % Arrows
  \draw[arr] (jira)     -- node[left, font=\tiny]{API REST} (etl);
  \draw[arr] (etl)      -- node[above, font=\tiny]{to\_sql()} (mysql);
  \draw[arr] (etl)      -- (streamlit);
  \draw[arr] (mysql)    -- node[right, font=\tiny]{localhost} (gateway);
  \draw[arr] (mysql)    -- node[above, font=\tiny]{connexion directe} (pbd);
  \draw[garr] (gateway) -- node[right, font=\tiny]{tunnel sécurisé} (pbs);

  % Network boundary
  \begin{pgfonlayer}{background}
    \node[draw=greenborder, dashed, rounded corners=6pt, fill=greenbg, fill opacity=0.1,
          fit=(etl)(mysql)(streamlit)(gateway),
          label={[font=\small\color{greenborder}]above:Réseau interne DXC}] {};
  \end{pgfonlayer}

\end{tikzpicture}
\end{center}

\vspace{0.4cm}

\subsection{Composants}

\begin{enumerate}[leftmargin=2em, itemsep=6pt]

  \item \textbf{Jira REST API}~: source de données officielle. Le pipeline
        Python se connecte via un token API personnel et récupère les tickets
        par pagination incrémentale. Seul le poste interne initie les requêtes
        sortantes.

  \item \textbf{Python ETL}~: collecte les données, applique toute la logique
        métier (burndown, matching des noms, calcul de l'engagement, affectation
        des squads) et alimente la base MySQL. Entièrement versionné sous Git.

  \item \textbf{MySQL 8.0 en conteneur Docker}~: base de données relationnelle
        légère, conteneurisée pour la reproductibilité. Un volume Docker persiste
        les données indépendamment du cycle de vie du conteneur. La directive
        \texttt{restart: always} assure le redémarrage automatique.

  \item \textbf{Streamlit}~: tableau de bord web interne accessible depuis le
        réseau local, permettant aux membres de l'équipe de consulter les données
        en temps quasi-réel sans accès à Power BI.

  \item \textbf{On-premises Data Gateway}~: passerelle Microsoft installée sur
        le poste hébergeant Docker. Elle crée un tunnel chiffré entre Power BI
        Service et la base MySQL locale, sans exposer de port public.

  \item \textbf{Power BI}~: couche de visualisation finale. Power BI Desktop
        se connecte directement en local~; Power BI Service se rafraîchit via
        le Gateway pour les rapports partagés.

\end{enumerate}

\vspace{0.3cm}

\subsection{Justification du choix}

\begin{greenbox}
\begin{itemize}[leftmargin=1.5em, itemsep=4pt]
  \item \textbf{Sécurité et conformité}~: les données ne quittent jamais le
        réseau DXC. Aucun prestataire cloud tiers ne stocke d'informations
        sensibles (noms d'employés, références clients, contenus de tickets).
        Conformité RGPD assurée par conception.
  \item \textbf{Coût nul}~: Docker, MySQL et le Gateway Microsoft sont
        entièrement gratuits. Aucun abonnement externe requis.
  \item \textbf{Performance}~: MySQL gère sans difficulté les 55\,000+ tickets
        actuels et peut absorber plusieurs années de données supplémentaires
        (estimation~: moins de 500 Mo pour 5 ans d'historique).
  \item \textbf{Reproductibilité}~: un simple \texttt{docker-compose up}
        reconstruit l'environnement complet sur n'importe quel poste DXC.
  \item \textbf{Évolutivité}~: si le volume de données croît significativement,
        le conteneur peut être migré vers un serveur dédié ou une VM interne
        sans modifier une seule ligne de code applicatif.
  \item \textbf{Maintenabilité}~: le code Python est versionné sous Git,
        testable, lisible par tout développeur. La logique métier n'est pas
        enfouie dans des formules DAX ou des requêtes M opaques.
\end{itemize}
\end{greenbox}

\vspace{0.4cm}

\subsection{Tableau comparatif récapitulatif}

\begin{longtable}{>{\small\bfseries}p{3.5cm}
                  >{\small\centering}p{2.5cm}
                  >{\small\centering}p{2.5cm}
                  >{\small\centering}p{2.5cm}
                  >{\small\centering\arraybackslash}p{2.5cm}}
  \toprule
  \textbf{Critère} &
  \textbf{Aiven (cloud)} &
  \textbf{Power BI natif} &
  \textbf{SharePoint} &
  \textbf{Docker local} \\
  \midrule
  \endhead
  \bottomrule
  \endfoot

  Sécurité données   & \textcolor{orangeborder}{Moyenne}  & \textcolor{orangeborder}{Moyenne}  & \textcolor{greenborder}{Bonne}    & \textcolor{greenborder}{\textbf{Excellente}} \\[3pt]
  Conformité RGPD    & \textcolor{redborder}{À valider}   & \textcolor{redborder}{À valider}   & \textcolor{greenborder}{Bonne}    & \textcolor{greenborder}{\textbf{Native}} \\[3pt]
  Coût               & \textcolor{greenborder}{Gratuit*}  & \textcolor{redborder}{Payant}      & \textcolor{greenborder}{Inclus}   & \textcolor{greenborder}{\textbf{Gratuit}} \\[3pt]
  Performance        & \textcolor{greenborder}{Bonne}     & \textcolor{redborder}{Limitée}     & \textcolor{redborder}{Mauvaise}   & \textcolor{greenborder}{\textbf{Bonne}} \\[3pt]
  Volume de données  & \textcolor{greenborder}{Illimité}  & \textcolor{redborder}{< 15k lignes}& \textcolor{redborder}{< 5k lignes}& \textcolor{greenborder}{\textbf{Illimité}} \\[3pt]
  Logique métier     & \textcolor{greenborder}{Python}    & \textcolor{redborder}{DAX/M}       & \textcolor{redborder}{Limitée}    & \textcolor{greenborder}{\textbf{Python}} \\[3pt]
  Maintenabilité     & \textcolor{greenborder}{Bonne}     & \textcolor{redborder}{Difficile}   & \textcolor{redborder}{Difficile}  & \textcolor{greenborder}{\textbf{Excellente}} \\[3pt]
  Disponibilité PBI  & \textcolor{greenborder}{24/7}      & \textcolor{greenborder}{24/7}      & \textcolor{greenborder}{24/7}     & \textcolor{orangeborder}{Poste allumé} \\[3pt]
  Dépendance externe & \textcolor{redborder}{Oui}         & \textcolor{redborder}{Oui}         & \textcolor{orangeborder}{Partielle}& \textcolor{greenborder}{\textbf{Aucune}} \\[3pt]

\end{longtable}

{\small * Tier gratuit limité, payant au-delà}

\vspace{0.5cm}
\noindent{\color{dxcpurple}\rule{\linewidth}{0.8pt}}
\begin{center}
  \small\color{greytext}
  Document interne — DXC Technology · Projet IRPAUTO
\end{center}

\end{document}
